% Chapter 8 -- translated by Claude (2026), continuing the voice and
% register of the Burkett-Hutcheson chapters.

\chapter{Symbolic simulation of alternating current circuits}

\section{General concepts}

Symbolic simulation in alternating current uses the same techniques
as symbolic simulation in direct current, with only two exceptions.

\section{The cSolve command}

To solve complex equations, the calculator has the \code{cSolve}
command, which works identically to the \code{solve} command. The
\code{cSolve} command is a powerful ally in obtaining numerical
answers from a symbolic simulation, when the equations involved are
complex. In other cases, when the equations used are real, the
\code{solve} command can be used.

\section{Complex unknowns}

To find the numerical value of a complex unknown, we must find two
numerical values: a real value and an imaginary value. Therefore, a
complex unknown is really two unknowns: one real and one imaginary.

Likewise, when we have a complex value as a known given value, we
really have two known values: a real value and an imaginary value.
In our first problem of this chapter, we will see how to find the
numerical value of two symbolic unknowns, using only one known
complex value.

\problem{057}

\emph{Problem: This problem has no figure. A circuit consists of a
resistance in series with a capacitance. Find their values, if when
applying a voltage of 240~V RMS at 200~Hz, the current taken from
the source is $1.2 + 1.6i$. Note that this current is not in RMS.}

Solution: First, we simulate the circuit according to the rigorous
treatment:

\begin{entry}
sq\ac("e,1,0,240.*√2;r,1,2,r;c,2,0,c",2*π*200.)
\end{entry}

Since there is only one element of each type, we have taken the
opportunity to name them \code{e}, \code{r} and \code{c}, which is a
licit and very elegant notation. If there were more than one element
of each type, they would have to be named more completely. We press
\key{ENTER}. After the status messages, `Done' appears. Next, we
order the calculator to find the values of \code{r} and \code{c}
that satisfy both the real part and the imaginary part of the
current that the problem gave us as a known value:

\begin{entry}
solve(real(ir)=1.2 and imag(ir)=1.6,{r,c})
\end{entry}

This command line is very interesting. Let us see what it means. We
have just run a simulation, using \code{r} and \code{c} as symbolic
values. Therefore, all the answers, including the current through
the elements, are expressed as functions of \code{r} and \code{c}.
These answers have a real part and an imaginary part. The current
that the circuit gave us as a known value also has a real part and
an imaginary part. What the command line is ordering the calculator
is: ``Use the \code{solve} command to find which values of \code{r}
and \code{c} satisfy the following two conditions: that the real
part of the current equal 1.2 amperes, and the imaginary part of the
current equal 1.6 amperes.'' We use the \code{solve} command and not
the \code{cSolve} command, because the two equations involved are
given in terms of real numbers. Why? The real part of a complex
number is a real number. The imaginary part of a complex number,
once it has been extracted with the \code{imag} command, is also a
real number. Therefore, the two equations given are composed of real
numbers, and the \code{solve} command can be used to solve them.

We press \key{ENTER}. The calculator gives us two pairs of answers.
One of these pairs has a negative value of \code{r}. The other pair
has a positive value of \code{r}. Obviously, the correct pair of
answers is the one with the positive value of \code{r}. These
values, rounded to three figures, are: \code{r}~$= 102.$ and
\code{c}~$= 5.86$E$-6$. The value of \code{r} is given in $\Omega$,
and that of \code{c} in F. If we want to obtain only the correct
answers, we indicate that $r > 0$, with this line:

\begin{entry}
solve(real(ir)=1.2 and imag(ir)=1.6,{r,c})|r>0
\end{entry}

\problem{058}

\emph{Problem: If $\omega = 500$~rad/s and $I_L = 2.5∠40°$, find
$v_s(t)$.}

Solution: We will use a symbolic simulation and then the
\code{solves} tool:

\begin{entry}
sq\ac("es,1,0,vs;r1,1,2,10;e2,2,3,(25∠-30°);r2,3,0,25;
      l1,3,0,.02",500):sq\solves()
\end{entry}

We press \key{ENTER}. After the status messages, `Done' appears.
Next, the tool executes. We enter \code{il1=(2.5∠40°)} as the
equation and \code{vs} as the unknown. We obtain the value of
\code{vs} on the screen, but it is in rectangular form and this way
it is not useful to us. We press \key{ENTER} to return to the home
screen. To see the value in polar form, we use
\code{sq\char`\\absang(vs)} in the entry line. We obtain
35.5∠58.9°. Therefore, the value of $v_s(t)$ is
$35.5\cos(500t + 58.9°)$~V.

\problem{059}

\emph{Problem: This problem has no figure. A phasor current of
$1∠0°$~A flows through the series combination of 1~$\Omega$, 1~H
and 1~F. At what frequency is the amplitude of the voltage across
the network twice the amplitude of the voltage across the
resistor?}

Solution: First, we simulate the circuit:

\begin{entry}
sq\ac("j,0,1,1;r,1,2,1;l,2,3,1;c,3,0,1",ω)
\end{entry}

Since there is only one element of each type, we have taken the
opportunity to name them with a single letter. We press \key{ENTER}.
After the status messages, `Done' appears. In this problem, it is
possible to find more than one answer, so the most adequate thing is
to use the \code{solve} command instead of the \code{solves} tool,
because the command will give us the multiple answers it finds,
while the tool will only give us one answer. So, we use the command:

\begin{entry}
solve(abs(v1)=2*abs(vr),ω)|ω>0
\end{entry}

Note that we suppressed the possibility of obtaining negative
frequencies, with the conditional at the end of the command line. We
will obtain two answers, and both are valid. If we press
\key{ENTER}, we obtain the answers in exact form:
$\omega = (\sqrt7 - \sqrt3)/2$ or $\omega = (\sqrt7 + \sqrt3)/2$. If
we press \key{$\diamond$} \key{ENTER}, we obtain the answers in
approximate form: $\omega = .457$ or $\omega = 2.19$.

\problem{060}

\emph{Problem: This problem has no figure. A 20~mH inductor and a
30~$\Omega$ resistor are in parallel. Find the frequency at which:
a) $\mathrm{abs}(Z_{in}) = 25~\Omega$, b)
$\mathrm{ang}(Z_{in}) = 25°$, c) $\mathrm{real}(Z_{in}) =
25~\Omega$, d) $\mathrm{imag}(Z_{in}) = 10~\Omega$.}

Solution: We use the \code{thevenin} tool:

\begin{entry}
sq\thevenin("r,1,0,30;l,1,0,.02",1,0)
\end{entry}

We press \key{ENTER}. Next, the tool executes. We select Passive and
AC, and press \key{ENTER}. We enter a symbolic frequency \code{w}
and press \key{ENTER}. After the status messages, \code{zth} appears
as a function of \code{w}. We press \key{ENTER} to return to the
home screen. To answer the questions, we use:

\begin{entry}
solve(abs(zth)=25,w)|w>0
\end{entry}

We obtain \code{w}~$= 2261$.

\begin{entry}
solve(angle(zth)=25°,w)|w>0
\end{entry}

We obtain \code{w}~$= 3217$.

\begin{entry}
solve(real(zth)=25,w)|w>0
\end{entry}

We obtain \code{w}~$= 3354$.

\begin{entry}
solve(imag(zth)=10,w)|w>0
\end{entry}

We obtain \code{w}~$= 3927.$ or \code{w}~$= 572.9$.

\problem{061}

\emph{Problem: In the network shown, find the frequency at which a)
$R_{in} = 550~\Omega$, b) $X_{in} = 50~\Omega$, c)
$G_{in} = 1.8$~mS, d) $B_{in} = -150~\mu$S. Remember that
$Z = R + iX$ and $Y = G + iB$.}

Solution: We use the \code{thevenin} tool:

\begin{entry}
sq\thevenin("r1,1,2,500;r2,2,0,100;l1,2,0,1E-3",1,0)
\end{entry}

We press \key{ENTER}. Next, the tool executes. We select Passive and
AC, and press \key{ENTER}. We enter a symbolic frequency \code{w}
and press \key{ENTER}. After the status messages, \code{zth} appears
as a function of \code{w}. We press \key{ENTER} to return to the
home screen. To answer the questions, we use:

\begin{entry}
solve(real(zth)=550,w)|w>0
\end{entry}

We obtain \code{w}~$= 100000$.

\begin{entry}
solve(imag(zth)=50,w)|w>0
\end{entry}

We obtain \code{w}~$= 100000$.

\begin{entry}
solve(real(1/zth)=.0018,w)|w>0
\end{entry}

We obtain \code{w}~$= 102062$.

\begin{entry}
solve(imag(1/zth)=-1.5E-4,w)|w>0
\end{entry}

We obtain \code{w}~$= 132953.$ or \code{w}~$= 52232$.

\problem{062}

\emph{Problem: Find $Y_{in}$ as a function of $\omega$ for the
network shown. Also find the resonant frequency $\omega_0$ and the
value of $Z_{in}$ at that resonant frequency.}

Solution: Theory teaches us that the resonant frequency is that
frequency at which the imaginary part of the input admittance equals
zero. As we know, the input admittance is the inverse of the input
impedance. We use the \code{thevenin} tool:

\begin{entry}
sq\thevenin("l,1,0,4.4E-3;r,0,1,1E4;c,1,2,1E-8;
            e,2,0,1E5*ir",1,0)
\end{entry}

We press \key{ENTER}. Next, the tool executes. We select Passive and
AC, and press \key{ENTER}. We enter a symbolic frequency \code{w}
and press \key{ENTER}. After the status messages, \code{zth} appears
as a function of \code{w}. We press \key{ENTER} to return to the
home screen. To answer the questions, we use:

\begin{entry}
solve(imag(1/zth)=0,w)|w>0
\end{entry}

With the command above these lines, we are asking the calculator:
``Using the \code{solve} command, tell me how much \code{w} is worth
when the imaginary part of the inverse of \code{zth} is zero,
considering \code{w} as positive.'' We press \key{ENTER}. We obtain
\code{w}~$= 45454.5$ as the resonant frequency.

\begin{entry}
zth|ans(1)
\end{entry}

With the command above these lines, we are asking the calculator:
``Tell me how much \code{zth} is worth when the immediately previous
answer holds.'' Placing \code{zth|ans(1)} is equivalent to placing
\code{zth|w=45454.5}. The \code{ans(\#)} command recalls the \#th
previous answer from the history area. We press \key{ENTER}. We
obtain 9999.9. This answer can be rounded to 10~k$\Omega$.

\problem{063}

\emph{Problem: For the network shown, find the resonant frequency
$\omega_0$ and the value of $Z_{in}$ at that frequency.}

Solution: We already know the technique:

\begin{entry}
sq\thevenin("r1,1,2,5;c,1,2,.01;r2,2,3,2;r3,3,0,5;l,3,0,.01",1,0)
\end{entry}

We press \key{ENTER}. Next, the tool executes. We select Passive and
AC, and press \key{ENTER}. We enter a symbolic frequency \code{w}
and press \key{ENTER}. After the status messages, \code{zth} appears
as a function of \code{w}. We press \key{ENTER} to return to the
home screen. To answer the questions, we use:

\begin{entry}
solve(imag(1/zth)=0,w)|w>0
\end{entry}

We press \key{ENTER}. We obtain \code{w}~$= 100.$ as the resonant
frequency.

\begin{entry}
zth|ans(1)
\end{entry}

We press \key{ENTER}. We obtain 2.38.
